%%%%%%%%%%%%%%%%%%%%%%%%%%%%%%%%%%%%%%%%%%%%%%%%%%%%%%%%%%%%%%%%%%%%%%
% Lớp 11 | Hàm số mũ và hàm số lôgarít | Phép tính lôgarít
% Dạng 1D6?2-4: Số e và bài toán lãi kép
% Số câu: 15  (sinh tự động bởi tach_cau.py; từ nay thầy sửa trực tiếp file này)
%%%%%%%%%%%%%%%%%%%%%%%%%%%%%%%%%%%%%%%%%%%%%%%%%%%%%%%%%%%%%%%%%%%%%%

%%%%%%%%%%%%%%%% TRẮC NGHIỆM 4 PHƯƠNG ÁN %%%%%%%%%%%%%%%%

% nguồn: Ok_TracNghiemLop11_Lan1
\begin{ex}%[1D6H2-4]%[ma:beab69cd]
	Một người gửi $100$ triệu đồng vào ngân hàng với lãi suất $0{,}4\% /$ tháng.  Biết rằng nếu không rút tiền ta khỏi ngân hàng thì cứ sau mỗi tháng, số tiền lãi sẽ được lập vào vốn ban đầu để tính lãi cho tháng tiếp theo. Hỏi sau $6$ tháng, người đó được lĩnh số tiền (cả vốn ban đầu và lãi) gần nhất với số tiền nào dưới đây, nếu trong khoảng thời gian này người đó không rút tiền ra và lãi xuất không thay đổi?
	\choice 
		 {$102\,160\,000$ đồng}
		{$102\,017\,000$  đồng}
		{\True  $102\,424\,000$ đồng}
		 {$102\,423\,000$ đồng}
	\loigiai{ 
		Ta có  $A_n=A_0\left(1+r\right)^n=100\,000\,000\left(1+\dfrac{0{,}4}{100}\right)^6=102\,424\,128$.
	 }
\end{ex}

% nguồn: Ok_TracNghiemLop11_Lan1
\begin{ex}%[1D6H2-4]%[ma:926e5004]
	Một điện thoại đang nạp pin, dung lượng pin nạp được tính theo công thức $Q(t)=Q_0\left(1-\mathrm{e}^{-t \sqrt{2}}\right)$, với $t$ là khoảng thời gian tính bằng giờ và $Q_0$ là đung lượng nạp tối đa (pin đầy). Thời gian nạp pin của điện thoại  tính từ lúc cạn pin cho đến khi điện thoại đạt được $90 \%$ dung lượng pintối đa là bao nhiêu? (kết quả làm tròn đến hàng phần trăm)
	 \choice
		 {$1{,}5$ giờ}
		{$1{,}61$ giờ} 
		 {\True $1{,}63$ giờ}
		{$1{,}65$  giờ}
	\loigiai{ 
		Ta có $\dfrac{Q(t)}{Q_0}=90\%\Leftrightarrow \left(1-\mathrm{e}^{-t \sqrt{2}}\right)=0{,}\Leftrightarrow  \mathrm{e}^{-t \sqrt{2}}=0{,}1\Leftrightarrow t=-\dfrac{\ln 0{,}1}{\sqrt{2}}\approx 1{,}63$.
	 }
\end{ex}


%%%%%%%%%%%%%%%% ĐÚNG/SAI %%%%%%%%%%%%%%%%

% nguồn: Ok_dungsai_lop11_Lan1
\begin{ex}%[1D6H2-4]%[ma:b3ccd8a6]
  Cho biểu thức $B=2\ln \sqrt{\mathrm{e}x}-\ln \dfrac{\mathrm{e}^2}{\sqrt{x}}+\ln 3 \cdot \log_3(\mathrm{e} x^2)$, với $x$ là số thực dương.
  \choiceTF
  {$B=\ln  x+\dfrac{3}{2}$}
   {\True Cho $\ln x =4$ thì $B=14$}
   {Cho $x =e^3$ thì $2B+3=\dfrac{15}{2}$}
  {\True Cho $x =6$ thì  $e^B=216\sqrt{6}$}
  \loigiai{
  Ta  có
  \begin{eqnarray*}
   B&=&2\ln \left( \mathrm{e} x\right)^{\frac{1}{2}}- \left(\ln\mathrm{e}^2-\ln x^{\frac{1}{2}}\right)+\ln  3 \cdot \dfrac{\ln (\mathrm{e}x^2) }{\ln 3} \\
   &= &\ln \left( \mathrm{e}  x\right)
   -\left( 2-\dfrac{1}{2}\ln x\right)+\ln  (\mathrm{e}x^2)\\
   &=&\left( \ln \mathrm{e}  +\ln x\right) -2+\dfrac{1}{2}\ln x +\ln e +\ln x^2\\
   &=&1+\ln x -2+\dfrac{1}{2}\ln x +1+2\ln  x\\
   &=&  \dfrac{7}{2}\ln x.
  \end{eqnarray*}
  \begin{itemchoice}
   \itemch Sai. Ta  có $B=\dfrac{7}{2}\ln x$.
    \itemch Đúng. Cho $\ln x =4$ thì $B=14$.
   \itemch Sai. Cho $x =e^3$  thì $2B+3=2\cdot\dfrac{7}{2}\ln(e^3)+3=2\cdot\dfrac{7}{2}\cdot 3+3=9$.
   \itemch Đúng. Cho  $x =6$ thì $B=\dfrac{7}{2}\ln 6$.\\
   $\Rightarrow e^B=e^{\tfrac{7}{2}\ln 6}=\left(e^{\ln 6}\right)^{\tfrac{7}{2}}=6^{\tfrac{7}{2}}=\sqrt{6^7}=6^3\cdot\sqrt{6}=216\sqrt{6}$. 
  \end{itemchoice}
 } 
\end{ex}


%%%%%%%%%%%%%%%% TRẢ LỜI NGẮN %%%%%%%%%%%%%%%%

% nguồn: Ok_TraLoiNgan_Lop11_Lan1
\begin{ex}%[1D6V2-4]%[ma:acc657de]
	Vào đầu mối tháng, ông An đều gửi yào ngân hàng  số tiền cố định 30 triệu đồng theo hình thức lãi kép với lãi suất $0{,}6\%$ /tháng. Tính tổng số tiền ông An có được sau tháng thứ hai. (đơn vị triệu đồng, làm tròn kết quả đến hàng phần mười).
	\shortans{60,5} 
	 \loigiai{
	Số tiền ông  An có được
	\begin{itemize}
		\item Sau tháng thứ nhất là $T_1=30+30 \cdot \dfrac{0{,}6}{100}=30\left(1+\dfrac{0{,}6}{100}\right)=30{,}18$ (triệu đồng). 
		\item Sau tháng thứ  hai là\\
		 $T_2=30+30\left(1+\dfrac{0{,}6}{100}\right)+\left[30+30\left(1+\dfrac{0{,}6}{100}\right)\right] \dfrac{0{,}6}{100}$
		 \\
		 $=\left[30+30\left(1+\dfrac{0{,}6}{100}\right)\right]\left(1+\dfrac{0{,}6}{100}\right)=30\left(1+\dfrac{0{,}6}{100}\right)+30\left(1+\dfrac{0{,}6}{100}\right)^2$
		 $\approx 60{,}5$ (triệu đồng)
	\end{itemize}
	 }
\end{ex}


%%%%%%%%%%%%%%%% TỰ LUẬN %%%%%%%%%%%%%%%%

% nguồn: Ok_TuanLuan_Toan11_Lan1
\begin{bt}%[1D6H2-4]%[ma:f9e2481e]
	Anh Hưng gửi tiết kiệm khoản tiền $700$ triệu đồng vào một ngân hàng với lãi suất $7 \%$/ năm theo hình thức lãi kép kì hạn $12$ tháng. Tính thời gian tối thiểu gửi tiết kiệm để anh Hưng thu được ít nhất $1$ tỉ đồng (cả vốn lẫn lãi). Cho biết công thức lãi kép là $T=A \cdot(1+r)^n$, trong đó $A$ là tiền vốn, $T$ là tiền vốn và lãi nhận được sau $n$ năm, $r$  là lãi suất/năm.
	\loigiai{ 
		 Ta có
		\begin{eqnarray*}
			T \geq 1000  \Leftrightarrow 700 \left( 1+7 \% \right)^n \geq 1000 &\Leftrightarrow& 1{,}07^n \geq \dfrac{10}{7} \\
			 &\Leftrightarrow& n \geq \log _{1,07}\left(\dfrac{10}{7}\right) \approx 5,27.
		\end{eqnarray*}
		Vậy thời gian gửi tiết kiệm phải ít nhất $6$ năm thì  anh Hưng mới thu được ít nhất $1$ tỉ đồng.
	} 
\end{bt}

% nguồn: Ok_TuanLuan_Toan11_Lan1
\begin{bt}%[1D6H2-4]%[ma:c8753858]
	Một khu rừng có trữ lượng gỗ là $4.10^5$m$^3$. Biết tốc độ sinh trưởng  của các cây lấy gỗ trong khu rừng này là $4\%$ mỗi năm. Hỏi sau 5 năm không khai thác, khu rừng sẽ có số mét khối gỗ là bao nhiêu?
	 \loigiai{
		 Nếu trữ lượng gỗ của khu rừng ban đầu là $ A $ thì sau năm thứ nhất, lượng gỗ có được là $ A+Ar=A(1+r) $ với $ r $ là tốc độ tăng trưởng mỗi năm.\\
		Sau năm thứ hai,  lượng gỗ có được là $A(1+r)+A(1+r)\cdot r=A(1+r)^2$.\\
		Theo phương pháp quy nạp, ta chứng minh được công thức tính lượng gỗ trong khu rừng là $ T_{n}=A(1+r)^n $  với $ A $ là lượng gỗ ban đầu, $ r $ là tốc độ tăng trưởng mỗi năm và $ n $ là số năm tăng trưởng của rừng.
		Vậy sau  5 năm, lượng gỗ trong khu rừng là:
		\[T_{5}=4\cdot 10^5\left(1+\dfrac{4}{100}\right)^5=486\,661{,}161\text{m}^3.\] 
	} 
\end{bt}

% nguồn: Ok_TuanLuan_Toan11_Lan1
\begin{bt}%[1D6H2-4]%[ma:bffe9c16]
	Giả sử số tiền gốc là $A$, lãi suất là $r \% /$ kì hạn gửi (có thể là tháng, quý hay năm) thì tồng số tiền nhận được cả gốc và lãi sau $n$ kì hạn gửi là $A(1+r)^n$. Bà  Hạnh gửi $100$ triệu vào tài khoản định kỳ tính lãi kép với lãi suất là $8 \% /$ năm. Tính số tiền lãi thu được sau $10$ năm.
	\loigiai{ 
		Áp dụng công thức tính lãi kép, sau $10$ năm số tiền cả gốc và lãi bà Hạnh thu về là $A(1+r)^n=100(1+0,08)^{10} \approx 215{,}892$  triệu đồng.\\
		 Suy ra số tiền lãi bà Hạnh thu về sau $10$ năm là $215{,}892-100=115{,}892$ triệu đồng.
	 }
\end{bt}

% nguồn: Ok_TuanLuan_Toan11_Lan1
\begin{bt}%[1D6H2-4]%[ma:5d5bb478]
	Anh Hưng gửi tiết kiệm khoản tiền $700$ triệu  đồng vào một ngân hàng với lãi suất $7\%$/ năm theo hình thức lãi kép kì hạn 12 tháng. Tính thời gian tối thiểu gửi tiết kiệm để anh Hưng thu được ít nhất 1 tỉ đồng (cả vốn lẫn lãi). Cho biết công thức lãi kép là $ T=A\cdot(1+r)^n $, trong đó $ A $ là tiền vốn, $ T $ là tiền vốn và lãi nhận được sau $ n $ năm, $ r $ là lãi suất/năm.
	 \loigiai{
		Ta có: $  T>1\,000\Leftrightarrow 700(1+7\%)^n\Leftrightarrow 1{,}07^n\ge \dfrac{10}{7}\Leftrightarrow n\ge \log_{1{,}07}\left(\dfrac{10}{7}\right)\approx5{,}27 $.\\
	Vậy thời gian gửi tiết kiệm phải ít nhất 6 năm thì anh Hưng mới thu được ít nhất  1 tỉ đồng.}
\end{bt}

% nguồn: Ok_TuanLuan_Toan11_Lan2
\begin{ex}%[1D6H2-4]%[ma:0abe3bef]
	Anh Hưng gửi tiết kiệm khoản tiền $700$ triệu đồng vào một ngân hàng với lãi suất $7 \%$/ năm theo hình thức lãi kép kì hạn $12$ tháng. Tính thời gian tối thiểu gửi tiết kiệm để anh Hưng thu được ít nhất $1$ tỉ đồng (cả vốn  lẫn lãi). Cho biết công thức lãi kép là $T=A \cdot(1+r)^n$, trong đó $A$ là tiền vốn, $T$ là tiền vốn và lãi nhận được sau $n$ năm, $r$ là lãi suất/năm.
	\loigiai{ 
		 Ta có
		\begin{eqnarray*}
			T \geq 1000 \Leftrightarrow 700 \left( 1+7 \% \right)^n \geq 1000  &\Leftrightarrow& 1{,}07^n \geq \dfrac{10}{7} \\
			&\Leftrightarrow& n \geq  \log _{1,07}\left(\dfrac{10}{7}\right) \approx 5,27.
		\end{eqnarray*}
		Vậy thời gian gửi tiết kiệm phải  ít nhất $6$ năm thì anh Hưng mới thu được ít nhất $1$ tỉ đồng.
	 }
\end{ex}

% nguồn: Ok_TuanLuan_Toan11_Lan2
\begin{ex}%[1D6H2-4]%[ma:e7cc8636]
	Một khu rừng có trữ lượng gỗ là $4.10^5$m$^3$. Biết tốc độ sinh trưởng của các cây lấy gỗ trong khu rừng này là $4\%$ mỗi năm. Hỏi sau 5 năm không khai thác,  khu rừng sẽ có số mét khối gỗ là bao nhiêu?
	 \loigiai{
		Nếu trữ lượng gỗ của khu rừng ban đầu là $ A $ thì sau năm thứ nhất, lượng gỗ có được là $ A+Ar=A(1+r) $  với $ r $ là tốc độ tăng trưởng mỗi năm.\\
		Sau năm thứ hai, lượng gỗ có được là $A(1+r)+A(1+r)\cdot r=A(1+r)^2$.\\ 
		Theo phương pháp quy nạp, ta chứng minh được công thức tính lượng gỗ trong khu rừng là $ T_{n}=A(1+r)^n $ với  $ A $ là lượng gỗ ban đầu, $ r $ là tốc độ tăng trưởng mỗi năm và $ n $ là số năm tăng trưởng của rừng.
		Vậy sau 5 năm, lượng gỗ trong khu rừng  là:
		\[T_{5}=4\cdot 10^5\left(1+\dfrac{4}{100}\right)^5=486\,661{,}161\text{m}^3.\] 
	 }
\end{ex}

% nguồn: Ok_TuanLuan_Toan11_Lan2
\begin{ex}%[1D6H2-4]%[ma:bff1170e]
	Giả sử số tiền gốc là $A$, lãi suất là $r \% /$ kì hạn gửi (có thể là tháng, quý hay năm) thì tồng số tiền nhận được cả gốc và lãi sau $n$  kì hạn gửi là $A(1+r)^n$. Bà Hạnh gửi $100$ triệu vào tài khoản định kỳ tính lãi kép với lãi suất là $8 \% /$ năm. Tính số tiền lãi thu được sau $10$ năm.
	 \loigiai{
		Áp  dụng công thức tính lãi kép, sau $10$ năm số tiền cả gốc và lãi bà Hạnh thu về là $A(1+r)^n=100(1+0,08)^{10} \approx 215{,}892$ triệu đồng.\\
		 Suy ra số tiền lãi bà Hạnh thu về sau $10$ năm là $215{,}892-100=115{,}892$ triệu đồng.
	} 
\end{ex}

% nguồn: Ok_TuanLuan_Toan11_Lan2
\begin{ex}%[1D6H2-4]%[ma:ef6b7d09]
	Anh Hưng gửi tiết kiệm khoản tiền $700$ triệu đồng vào một ngân hàng với lãi suất $7\%$/ năm theo hình thức lãi kép kì hạn 12 tháng. Tính thời gian tối thiểu gửi tiết kiệm để anh Hưng thu được ít nhất 1 tỉ  đồng (cả vốn lẫn lãi). Cho biết công thức lãi kép là $ T=A\cdot(1+r)^n $, trong đó $ A $ là tiền vốn, $ T $ là tiền vốn và lãi nhận được sau $ n $ năm, $ r $ là lãi suất/năm.
	 \loigiai{
		Ta có: $ T>1\,000\Leftrightarrow 700(1+7\%)^n\Leftrightarrow 1{,}07^n\ge \dfrac{10}{7}\Leftrightarrow n\ge  \log_{1{,}07}\left(\dfrac{10}{7}\right)\approx5{,}27 $.\\
	Vậy thời gian gửi tiết kiệm phải ít nhất 6 năm thì anh Hưng  mới thu được ít nhất 1 tỉ đồng.}
\end{ex}

% nguồn: Ok_TuanLuan_Toan11_Lan2
\begin{ex}%[1D6H2-4]%[ma:d0fdfcb2]
	Cường độ một trận động đất $M$ được cho bởi công thức $M=\log A-\log A_0$, với $A$ là biên độ rung chấn tối đa và $A_0$ là một biên độ chuẩn (hằng số). Đầu thế kỷ $20$, một trận động đất ở San Francisco có cường độ $8{,}3$ độ Richter. Trong cùng năm đó, trận động đất khác ở gần đó đo được $7{,}1$ độ  Richter. Hỏi trận động đất ở San Francisco có biên độ gấp bao nhiêu trận động đất này
	 \loigiai{
		Xét trận động đất ở San Francisco, ta có 
		$$8{,}3=\log A-\log A_0\Leftrightarrow  8{,}3=\log \dfrac{A}{A_0}\Leftrightarrow \dfrac{A}{A_0}=10^{8{,}3}\Leftrightarrow A=A_0\cdot 10^{8{,}3}.$$
		Xét  trận động đất khác gần đó, ta có\\
		$$7{,}1=\log A-\log A_0\Leftrightarrow 7{,}1=\log \dfrac{A}{A_0}\Leftrightarrow \dfrac{A}{A_0}=10^{7{,}1}\Leftrightarrow A=A_0\cdot  10^{7{,}1}.$$
		Từ đó ta tính tỉ lệ $\dfrac{A_0\cdot 10^{8{,}3}}{A_0\cdot  10^{7{,}1}}=10^{1{,}2}\approx 15{,}8$.\\
	Vậy trận động đất ở San Francisco có biên độ gấp $15{,}8$ trận động đất  khác gần đó.}
\end{ex}

% nguồn: Ok_TuanLuan_Toan11_Lan2
\begin{ex}%[1D6H2-4]%[ma:b38243af]
	Một học sinh $A$ khi đủ $18$ tuổi được cha mẹ cho $200\,000\,000$VNĐ. Số tiền này được bảo quản trong ngân hàng MSB với kì hạn thanh toán $1$ năm và học sinh $A$ chỉ nhận được số tiền này khi học xong $4$ năm đại học. Biết rằng khi đủ $22$ tuổi, số tiền mà học sinh $A$ được nhận  sẽ là $243\,101\,250$ VNĐ. Vậy lãi suất kì hạn một năm của ngân hàng MSB là bao nhiêu?
	 \loigiai{
		Gọi $P$ là số tiền gốc, $r$ là lãi suất một năm, $n$ là số năm gửi và $S_n$ là số tiền  nhận được sau $n$ năm.\\
		Theo đề bài ta có $P = 200\,000\,000$; $S_4 = 243\,101\,250$ và $n  = 22 - 18 = 4$.\\
		Áp dụng công thức lãi kép $S_n=P\cdot(1+r)^n$,\\ 
		 ta có phương trình
		$243\,101\,250  = 200\,000\,000 \cdot (1+r)^4$.\\
		Suy ra $(1+r)^4 = \dfrac{243\,101\,250}{200\,000\,000} = 1{,}21550625$.\\ 
		Do đó $1+r = \sqrt[4]{1{,}21550625}=1{,}05  \Rightarrow r = 0{,}05 = 5\%$.\\
		Vậy lãi suất kì hạn một năm của ngân hàng MSB là $5\%$. 
	 }
\end{ex}

% nguồn: Ok_TuanLuan_Toan11_Lan2
\begin{ex}%[1D6H2-4]%[ma:91e31ba1]
	Người ta sử dụng công thức $S=A\cdot e^{n\cdot r}$ để dự báo dân số của một quốc gia, trong đó $A$ là số dân của năm lấy làm mốc tính, $S$ là số dân sau $n$ năm và $r$ là tỉ lệ gia tăng dân số hàng năm.  Biết rằng năm $2\,001$ dân số của Việt Nam là $78\, 685\, 800$ người. Giả sử tỉ lệ tăng dân số hàng năm không đổi là $1{,}2\%$, hỏi dân số nước ta đạt $110$ triệu người vào năm nào?
	 \loigiai{
		Theo công thức  tăng trưởng mũ $S=A\cdot{e^{n\cdot r}}$.\\
		Suy ra $110\,000\,000=78\,685\,800\cdot e^{1{,}2\%\cdot n}\Leftrightarrow n=\dfrac{1}{1{,}2\%}\ln \dfrac{110\,000\,000}{78\,685\,800}\approx  27{,}91$.\\
		Nên sau $28$ năm thì dân  số Việt Nam năm nào sau đây đạt $110$ triệu người.\\
		Vậy dân số đạt  $110$ triệu người vào năm $2029$.
	 }
\end{ex}
